\documentclass{report}
\usepackage[utf8]{inputenc}
\usepackage[english]{babel}
\usepackage[]{graphicx}
\graphicspath{ {Figures/} }
\usepackage{multirow}
\usepackage{enumitem}
\usepackage{tabularx}
\usepackage{booktabs}
\usepackage{ltablex}
\usepackage{amsmath}
\usepackage[parfill]{parskip}
% Set page margins
\usepackage[top=100pt,bottom=100pt,left=68pt,right=66pt]{geometry}
% Changes the style of chapter headings
\usepackage{titlesec}
\titleformat{\chapter}
   {\normalfont\LARGE\bfseries}{\thechapter.}{1em}{}
% Change distance between chapter header and text
\titlespacing{\chapter}{0pt}{50pt}{2\baselineskip}
\newcommand{\tabincell}[2]{\begin{tabular}{@{}#1@{}}#2\end{tabular}} 

\pagenumbering{roman}

\begin{document}
% Title
\begin{titlepage}
	\clearpage\thispagestyle{empty}
	\centering
	\vspace{1cm}

	% Titles
	% Information about the University
	{\normalsize The University of Melbourne \\ 
		School of Computing and Information Systems \\
		SWEN90016 Software Processes and Management\\
		Semester 1 – 2020 \par}
		\vspace{3cm}
	{\Huge \textbf{JJFresh -- Project Management Plan}} \\
	%\vspace{1cm}
	%{\large \textbf{xxxxx} \par}
	\vspace{3.5cm}
	{\normalsize Zhaofeng Qiu 1101584\\ % \\ specifies a new line
	             Chongjing Zhang 1055520\\
	             Pin Wang 1056745 \\
	             Yicun Tian 1088217 \\
	             Hongkang Li 977063\par}
	\vspace{3cm}
    
    \centering \includegraphics[scale=0.12]{logo.pdf}
    
    \vspace{0.5cm}
		
	% Set the date
	{\normalsize \today \par}
	\pagebreak
\end{titlepage}

\chapter*{Executive Summary}\label{sec:ESummary}
Jess and James plans to create an online store of fruits and vegetables as a JJFresh online sales platform for users to select and place orders. In this way, they want to maintain their existing customer base and remain competitive with other stores. In view of the insufficient budget of Jess and James, the students of registered software project management of the University of Melbourne are employed to develop this system. Based on the consideration of the changes in the needs of the project itself and the size of the development team, they chose an agile Scrum model, which will be assigned tasks to members by a short Sprint. They plan to deliver a product that meets core needs to Jess and James within 2 weeks, and complete 3 iterations of development within 4 weeks to deliver a complete website. The product will implement an online sales platform for Jess and James, thereby bringing them additional economic benefits. On the other hand, for the development team, the project will accumulate actual development experience, which will improve their competitiveness in future. 

\addcontentsline{toc}{chapter}{Executive Summary}
\pagebreak
\clearpage

\tableofcontents
\pagebreak

\chapter{Introduction}
\pagenumbering{arabic}
\section{Purpose of document}
   This document details how the development team plans, implements, and monitors the development of the JJFresh online store. According to the project requirements specification, this document describes in detail the scope of the project and possible risks, and how to control it. The development team also described the technologies that will be used throughout the development process. In addition, Scrum Master made a project plan for the development team's size and project needs. This document will serve as one of the true sources of all parties involved in the entire development cycle, such as Scrum Master, Product Owner and developers.

\section{Audience of document}
   All five people in the team participated in the preparation of the document. Firstly, Yicun Tian wrote the Executive Summary (part 2), Introduction (part 4), Roles and Responsibilities (part 6.1), Communication Plan (part 6.2) and Risk management (part 6.3) of the document.Secondly,Hongzhang Li completed Project Information (part 5). Moreover, Pin Wang and Chongjing Zhang as developers completed Technology (part 6.4). Zhangfeng Qiu, who served as Scrum Master, completed the Project Planning (6.5) and final proofreading of all document.

\section{Evolution of document}

\chapter{Project Information}
%%%%%%%%%%%%%%%%%%%%%%%%%%%%%%%%%%%%%%%%%%%%%%%%%%%%
\section{Key Stakeholders}
\begin{enumerate}
  \item Jess and James: they are product owner who define the features, contents of this project, JJFresh online store website, and responsible for the benefits of it.
  \item Team: In this project, a group of students will come to develop it under the employment of the product owner. The project cannot be developed if there no people come to implement it.
  \item Customer: Consumption by customers is the whole of the business in this project. Customers will purchase fruit and vegetable on Online store, and the consumption is the majority of the income of JJFresh.
\end{enumerate}
%%%%%%%%%%%%%%%%%%%%%%%%%%%%%%%%%%%%%%%%%%%%%%%%%%%%%

\section{Scope}
\subsection{What is in-scope?}
\begin{enumerate}
  \item Project goal: establishing an online organic fruit and vegetable store.
  \item Project function: purchase function, administration manage order function and users information management.
  \item Project time: 1 month.
  \item Project deliverable: the website version can provide the function of normal business.
\end{enumerate}
\subsection{What is out-of-scope?}
\begin{enumerate}
  \item Customers: This project does not contain deal with the effect of attracting customers. How many customers will visit and consumption in the online store has nothing to do with the development.
  \item The sales of the online store: When the project is finished, no matter what sales of the online store, it has nothing to do with the project development. Jess and Jame should be responsible for the daily operation of JJFresh online store, including make some discount to attract more people.
  \item Fruit and vegetable suppliers: Because of the online store can sell more products than physical store, so the more suppliers need to be chosen. However, Suppliers have a direct relationship with product owner. This project does not contain deal with the supplier of the online store. Jess and Jame should concern about choosing which suppliers and have the business with them.
\end{enumerate}
%%%%%%%%%%%%%%%%%%%%%%%%%%%%%%%%%%%%%%%%%%%%%%%%%%%%%
\section{Delivery approach / SDLC - Agile}
LHK

Agile Scrum is the popular model that be used for website project.

This project is developed through a student group. Agile Scrum team are typically 6-9 people, without any traditional software engineers, and the team members work together to complete the set of work they have collectively committed to complete within a sprint (1). Compared to the other models, for instance, Incremental SDLC requires more resources and management attention, but the group members are students who need to spend more time on other subjects and cannot obtain enough resources like the professional software engineer.

This project is the online store website which greatly affected by market change and there is may be some change after the process in the underway. For example, product selection in the project maybe add more size of box or more type of box because of the market change during the development. Waterfall SDLC is difficult to accommodate these changes in the development, while Agile Scrum can adapt regular in the changing circumstances.

It is necessary to release the website as soon as possible to promote the online store. The first version can have only the product list and the introduction of online store to attract customers. When the final build online, JJFresh will have many customers who cannot wait to try its products. Agile Scrum can satisfy this point. In the development of this project, Agile Scrum can deliver many versions of project, and face to meet the increasing customer demands and needs. Compare to the Incremental Model, both of them delivery many build.

We choose to use Agile SDLC out of Formal, Agile and Hybrid, because we prefer the following advantages of agile compares to the limitations of a formal model such as waterfall.

Advantages of the Agile Model:
\begin{itemize}
  \item It is focused client process. So, it makes sure that the client is continuously involved during every stage. And it can reduce waste and rework and increase the satisfaction of clients.
  \item Agile teams are extremely motivated and self-organized so it likely to provide a better result from the development projects.
  \item Agile software development method assures that quality of the development is maintained
  \item The process is completely based on the incremental progress. Therefore, the client and team know exactly what is complete and what is not. This reduces risk in the development process.
  \item Agile intends small or mid-sized dedicated teams with high coordination, which is suitable for our team with 5 members.
\end{itemize}
Limitations of Waterfall Model:
\begin{itemize}
  \item Very difficult to move back to makes changes in the previous phases. It's unknown whether the client will change their requirements.
  \item The testing process starts once development is over. Hence, it has high chances of bugs to be found later in development where they are expensive to fix.
  \item Waterfall involves large teams. It decreases coordination among team members. While we have only 5 people in the team.
\end{itemize}
%%%%%%%%%%%%%%%%%%%%%%%%%%%%%%%%%%%%%%%%%%%%%%%%%%%%%
\section{Business Value (Financial \& Non-Financial Benefits)}
%\begin{itemize}
%  \item JJFresh would gain more users due to the extension of business mode, so Jess and James can get more economic %profits from the increased number of orders delivered by increased users.
%  \item Customers would be able to buy fruits and vegetables in a more convenient way with quality as good as usual.
%  \item The suppliers can get more orders from JJFresh, and also get more economic profit.
%\end{itemize}
\subsection{Financial benefit}
For Jess and Jame, online store can save the cost of hiring clerks, management and consumption of water and electricity in physical store, while it can bring the potential considerable income and alleviate the pressure caused by poor sales of the physical store.

For student in the team, they may have the reward after when the project is finished.

For customers, the product of online store often lower than the physical store.

\subsection{Non-financial benefit}
For Jess and Jame, the risk of establishing the online store is less than physical store. For instance, the cost of establishing and maintaining the website is less than physical store. Compare to physical store, online store can reduce the risk of products spoilage because of cargo backlog. Online store can serve more customers who may live far away from JJFresh. Hence, Jess and James may have more orders than the physical store.

For student in the team, they will obtain their experience in the development of the project. On one hand, the experience of development may attract students toward world of software project development and may help him/her in the career path on it. On the other hand, the experience can improve students' teamwork ability, which is important in the future.

For customers, online store provides the convenient platform to purchase vegetable and fruit. they can consumption product at night without worrying about the dangerous of going to the physical store. Delivery service can not only also saves the time, but send the products to the customers at the special period. For example, government ask people to reduce unnecessary travel because of coronavirus recently, delivery service will help people to get their products safety.

xxxlhk

JJFresh is currently busy in the morning but almost empty in the afternoon. The deployment of this project will keep JJFrsh's peak business hours to retain old customers who shop in the morning, and shift the afternoon business hours from offline to online, which offers delivery service for current and potential customers who cannot go to JJFresh store at that time. The online store will operate 24 / 7 which is much longer compare to the old business hours, thus provides customers with the convenience of placing orders at any time. With these advantages, the online store project would attack more customers for JJFresh. Therefore can increase fruit and vegetable orders, Jess and James can gain financial benifits, and avoid JJFresh from being closed.

The Deployment of the online store website can also increase the chance of JJFresh being searched by a search engine of a browser, which can improve the exposure, thereby improving brand awareness.

If this business mode achieves success, JJFresh can lead other retail stores to transition from offline to online.
%%%%%%%%%%%%%%%%%%%%%%%%%%%%%%%%%%%%%%%%%%%%%%%%%%%%%
\section{Constraints}
\begin{itemize}
  \item team members lack of experience in agile
  \item team members are unfamiliar with wix/front end techniques/other techniques?
  \item wix isn't compatible with other front end languages/frameworks, cannot be customized if needed?
  \item JJFresh didn't do market analysis, whether this business mode will success is unknown
  \item time is limited to 1 month
\end{itemize}
%%%%%%%%%%%%%%%%%%%%%%%%%%%%%%%%%%%%%%%%%%%%%%%%%%%%%

\chapter{Project Governance}
\section{Roles and Responsibilities}
\section{Communication Plan}
   Our group first uses “slack” for the main information transfer, such as publishing message notifications, file sharing, etc. Secondly, using “trello” to arrange the various stages of the project plan and the content is managed by the scrum master. Furthermore, “github” is used to manage output products, such as the release of code and report.
   According to the scrum development cycle and the difficulty of this project development, we will conduct a 15-minute meeting discussion on zoom every day, mainly to summarize the results of yesterday and the arrangement of today's tasks. If a team member is absent from a meeting, the person in charge will use email or telephone to communicate with him.

\section{Risk Management}
\section{Technology}
\section{Project Planning}
The SDLC of our project is SCRUM. Since the key requirements for the initial development provided by Jess and James are fixed, a Fixed-Scope Release Planning would be used to plan our project.

\subsection{Product Backlog}
The detail of product backlog is shown in table \ref{tab:productBacklog}. BFTB Score is the abbreviate of the Bang-for-the-Buck Score, which is a way of measuring how to get the most value in the shortest time. It is used to help assess the priority of the features. The BFTB Score is calculated by the formula below:
$$
  \text{BFTB Score} = \frac{\text{Value Point}}{\text{Store Point}}
$$

The story points of different features are estimated by the development team and the SCRUM master, based on the volume, risk, uncertainty and complexity of the features. The value points of them are estimated by the product owners. \textbf{Fabonacci sequence} is used to provided relative estimation to both story point and value point. Each story point equals about an hour of work.

Except for the tree future user stories at the end of the product backlog table, all others are must-have stories because they are key requirements defined by Jess and James.
\clearpage
\begin{tabularx}{\linewidth}{%
  >{\raggedright\arraybackslash}p{1.5cm}%
  >{\raggedright\arraybackslash}X%
  >{\raggedright\arraybackslash}X%
  p{1cm}p{1cm}p{1cm}
  }
  \toprule
  Feature & User Story & Tasks & Story Point & Value Point & BFTB Score\\
  \midrule
  Browse product menu 
  & As a customer, I want a menu that shows all the products, so that I can know what products are the website selling. 
  & 1.Provide a page for showing the list of the products\textit{(3-hour)}; 2.Show the types of the products\textit{(1-hour)}; 3.Show the price of the products\textit{(1-hour)}; 4.The price would change automatically\textit{(1-hour)}; 5.Show the pictures of products\textit{(1-hour)}.
  & 8
  & 5
  & 0.625\\
  \midrule
  Customer sign up
  & As customers, I want to sign up for the website, so that I can a member of the website.
  & 1.Check whether the information of the user is valid\textit{(1-hour)}; 2.Add the user to the database and send a confirming email to the new member if valid\textit{(1-hour)}.
  & 2
  & 2
  & 1
  \\
  \midrule
  Customer sign in and sign out
  & As customers, I should be able to sign in and sign out, so that I can manage my account and orders.
  & 1.Sign in, check whether the user exist and whether the user's password correct\textit{(1-hour)}; 2.Sign out\textit{(1-hour)}; 3.Handle the problem of forgetting the user password (Allow password reset)\textit{(1-hour)}.
  & 3
  & 3
  & 1
  \\
  \midrule
  Customer add and edit client information
  & As customers, after login, I can add or edit my information like home address, user name, email address, and contact number.
  & 1.Add and edit the name, email address, home address\textit{(1-hour)}; 2.Add and edit up to three multiple contact phone numbers\textit{(1-hour)}.
  & 2
  & 3
  & 1.5
  \\
  \midrule
  Add products to shopping cart
  & As customers, after login, I can add or edit my information like home address, user name, email address, and contact number.
  & 1.Select size, type and amount of the products and add them to shopping cart.\textit{(2-hour)}
  & 2
  & 5
  & 2.5
  \\
  \midrule
  Manage shopping cart
  & As a customer, I want to manage my shopping cart, so that I can remove those I don't want or add some more products.
  & 1.Select the items to pay\textit{(1-hour)}; 2.Change the amount of items\textit{(1-hour)}; 3.Delete items\textit{(1-hour)}; 4.View the price of the items selected\textit{(1-hour)}; 5.Provide a page to show all the items in the shopping cart\textit{(3-hour)}.
  & 8
  & 5
  & 0.625
  \\
  \midrule
  Check out shopping cart
  & As a customer, I want to choose the day and time for delivery, so that I can get the fruit on the right day.
  & 1.Choose a day and time for delivery.\textit{(1-hour)}. Only two bookings are allowed in a particular-hour.); 2.Send a confirming email to the customer\textit{(1-hour)}; 3.Show the day and time which is valid\textit{(2-hour)}; 4.Add an order to the user's order list\textit{(1-hour)}; 5.Handle synchronization problem\textit{(2-hour)}.
  & 8
  & 5
  & 0.625
  \\
  \midrule
  Cancel orders
  & As a customer, I need to get the ability to cancel an order so that I can modify an order and then re-create.
  & 1.Cancel order\textit{(1-hour)}; 2.Send emails to both customer and admin\textit{(1-hour)}.
  & 2
  & 2
  & 1
  \\
  \midrule
  Build database
  & As an admin, I need a database, so that I can store all the information of customers orders.
  & 1.Define the order model\textit{(1-hour)}; 2.Define the user information model\textit{(1-hour)}; 3.Define the admin information model\textit{(1-hour)}; 4.Create table and relation\textit{(2-hour)}; 5.Test\textit{(1-hour)}
  & 8
  & 5
  & 0.625
  \\
  \midrule
  Admin sign in and sign out
  & As an admin, I want to provide my username and password, so that I can register an admin account.
  & 1.Sign in\textit{(1-hour)}; 2. Sign out\textit{(1-hour)}.
  & 2
  & 2
  & 1
  \\
  \midrule
  View and manage orders
  & As an admin, I want the specific information of the user's order, so that I can packaging user orders.
  & 1.Confirm order\textit{(1-hour)}; 2.Order rank by date\textit{(1-hour)}; 3.Cancel order and provide reason\textit{(1-hour)}; 4.Provide a page to show order list with order information attach to it\textit{(4-hour)};
  & 8
  & 5
  & 0.625
  \\
  \midrule
  Manage product infomation
  & As an admin, I want to manage product information, so that I can change the price or picture of my products.
  & 1.Change the price of products\textit{(1-hour)}; 2.Change the pictures of products\textit{(1-hour)}; 3.Change the type of products\textit{(1-hour)}; 4.Add products\textit{(1-hour)}; 5.Delete products\textit{(1-hour)}.
  & 5
  & 3
  & 0.6
  \\
  \midrule
  Register admin account
  & As an admin, I need to get the ability to register for a new admin account, so that I can get more admin accounts if I need them.
  & 1.Provide way for admin account registration\textit{(1-hour)}.
  & 1
  & 1
  & 1
  \\
  \midrule
  Edit admin account information
  & As an admin, I may want to edit my admin account information, so that I can change the password.
  & 1.Provide way for admin account information editing\textit{(1-hour)}.
  & 1
  & 1
  & 1
  \\
  \midrule
  Add multiple bookings for the same slot
  & As an admin, I want to add multiple bookings for the same slot in the future. So that I  can employ others to deliver boxes.
  & 1.Allow multiple bookings for the same slot\textit{(6-hour)}.
  & 8
  & 2
  & 0.25
  \\
  \midrule
  Using AI to calculate delivery routes and times
  & As an admin, I want to use AI technology to help me calculate an optimal route and time for delivery.
  & 1.Create a model and an algorithm for the AI to learn and help calculate the optimal routes and time for delivery\textit{(12-hour)}.
  & 13
  & 3
  & 0.231
  \\
  \midrule
  Extended to allow for payment online.
  & As an admin, I want to be paid online, so that I can save my time and improve the delivery efficiency.
  & 1.Provide ways for online payment\textit{(5-hour)}; 2.Secure user information\textit{(8-hour)}.
  & 13
  & 3
  & 0.231
  \\
  \bottomrule
  \\
  \caption{Product Backlog}  
  \label{tab:productBacklog}
\end{tabularx}  

\subsection{Must-have Story Points}
Without calculating the features needed in the future, the total number of must-have story points is 57.
$$
SP_{total} = 57 \text{ (Story Point)}
$$

\subsection{Velocity Estimating}
There are two developers in our team. Since both of them are students, each of them can only spend about 1.2 to 1.5 hours a day on the project. They can work five days a week, so the total working hours of a week are 12 to 15 hours. We expect each story point to be equivalent to one hour of working time. So the min velocity ($V_{min}$) for a two-week sprint is about 24 story points, and the max velocity ($V_{max}$) is about 30 story points. And by using the formula shown below, we can get our minimum and maximum number of sprints, which are 2 and 2.5, respectively.
$$
S_{min} = \frac{SP_{total}}{V_{max}}
\text{, } 
S_{max} = \frac{SP_{total}}{V_{min}}
$$

\subsection{Sprint Planning}
\subsubsection{Sprint Cycle}
\label{sec:sprintCycle}
The development phase of our project would start on \textbf{Monday, April 25}. Since the final delivery due date of the whole project is \textbf{June 1}, we decided to divide our development into three phases, each corresponding to a sprint.

The first phase is to let the website go online without the function of purchase. The website can display all kinds of information about the products and provide the service of user registration. At the same time, the database is initially established to record the data of users and products. The first sprint will last for two weeks, \textbf{from Monday, April 27 to Friday, May 8}. 


In the second phase, we need to realize the function that users place orders to purchase products, and sellers can view and modify user orders. Preliminary testing will also take place at this phase. The second sprint will last for two weeks, \textbf{from Monday, May 11 to Friday, may 22}.

The third stage is mainly based on user feedback to repair the vulnerability, carry out more tests and improve the existing functions. If some functions left over from the first two sprints are not implemented, they will also be solved in this sprint. As the project is nearing its end and there is less work left, the sprint in this phase will last only one week, \textbf{from Monday, May 25 to Friday, May 29.}

On \textbf{the first Monday of every sprint}, a Sprint Planning meeting would be held to select high priority items from the product backlog that the development team can commit to delivering in a single Sprint. The select itmes would be add to Sprint Backlog.

On \textbf{the last Friday of each sprint}, a Sprint Review meeting would be held to demonstrate the new features to Stakeholders, and a Sprint Retrospective meeting would be held to review the sprint.

A 15-minute Daily SCRUM would be held \textbf{everyday} to start up the jobs. Everyone has to share what he/she did yesterday, what he/she plan to do, and what obstacles are slowing him/her.

\subsection{The First Sprint Plan}
As mention in section \ref{sec:sprintCycle}, the main goal of the first sprint is to launch a webpage where customers can browse products. Considering both the BFTB Score, which is shown in table \ref{tab:productBacklog} and the actual development needs, the Sprint Backlog of the first sprint is shown in table \ref{tab:firstSprintBacklog}. The total velocity and the delivery value of the first sprint is 30 story points and 25 value points, respectively. 
\\
\begin{tabularx}{\linewidth}{%
  l%
  >{\raggedright\arraybackslash}p{2cm}%
  >{\raggedright\arraybackslash}X%
  >{\raggedright\arraybackslash}p{1cm}}
  \toprule
  Index & Feature & Tasks & Story Point\\
  \midrule
  1 
  & Build database 
  & 1.Define the order model\textit{(1-hour)}; 2.Define the user information model\textit{(1-hour)}; 3.Define the admin information model\textit{(1-hour)}; 4.Create table and relation\textit{(2-hour)}; 5.Test\textit{(1-hour)}
  & 8
  \\
  \midrule
  2 
  & Browse product menu
  & 1.Provide a page for showing the list of the products\textit{(3-hour)}; 2.Show the types of the products\textit{(1-hour)}; 3.Show the price of the products\textit{(1-hour)}; 4.The price would change automatically\textit{(1-hour)}; 5.Show the pictures of products\textit{(1-hour)}.
  & 8
  \\
  \midrule
  3
  & Customer sign up
  & 1.Check whether the information of the user is valid\textit{(1-hour)}; 2.Add the user to the database and send a confirming email to the new member if valid\textit{(1-hour)}.
  & 2
  \\
  \midrule
  4
  & Customer sign in and sign out
  & 1.Sign in, check whether the user exist and whether the user's password correct\textit{(1-hour)}; 2.Sign out\textit{(1-hour)}; 3.Handle the problem of forgetting the user password (Allow password reset)\textit{(1-hour)}.
  & 3
  \\
  \midrule
  5
  & Customer add and edit client information
  & 1.Add and edit the name, email address, home address\textit{(1-hour)}; 2.Add and edit up to three multiple contact phone numbers\textit{(1-hour)}.
  & 2
  \\
  \midrule
  6
  & Manage product infomation
  & 1.Change the price of products\textit{(1-hour)}; 2.Change the pictures of products\textit{(1-hour)}; 3.Change the type of products\textit{(1-hour)}; 4.Add products\textit{(1-hour)}; 5.Delete products\textit{(1-hour)}.
  & 5
  \\
  \midrule
  7
  & Admin sign in and sign out
  & 1.Sign in\textit{(1-hour)}; 2. Sign out\textit{(1-hour)}.
  & 2
  \\
  \bottomrule
  \\
  \caption{The first Sprint Backlog (4.27 - 5.8)}  
  \label{tab:firstSprintBacklog}
\end{tabularx}
\clearpage

Trello is used as the PMP board to help us manage our SCRUM process. After the first sprint planning, the PMP board of our project is shown in Figure \ref{fig:firstSprintPmpBoard}. The label of the card is about the value point of the feature. The Story point of the feature is shown in the upper right corner of the card.

\begin{figure}[htp]
\centering
\includegraphics[width=\textwidth]{Figures/pmpBoardFirstSprint.png}
\caption{The PMP Board set after the first Sprint Planning}
\label{fig:firstSprintPmpBoard}
\end{figure}


\chapter{Project Execution, Monitoring and Control}
\section{Project Status: Friday Week 9}
\subsection{Process Related Artefacts}
\subsection{Product Related Artefacts}
\subsection{Risk Monitoring and Control}
\section{Project Status: Friday week 10}
\subsection{Process Related Artefacts}
\subsection{Product Related Artefacts}
\subsection{Risk Monitoring and Control}
\section{Project Status: Friday week 10}
\subsection{Process Related Artefacts}
\subsection{Product Related Artefacts}
\subsection{Risk Monitoring and Control}



\end{document}